\chapter{概率论的基础知识}

\intro{
	概率论是研究随机现象数量规律的数学分支，也是数理统计、随机过程等学科的理论基础。
}

\section{随机试验与样本空间}

\begin{mydefinition}{随机试验}
	对随机现象的观察，称为\textbf{随机试验}（简称试验），记为 $E$。
\end{mydefinition}

随机试验具有以下三个特点：
\begin{enumerate}
	\item 可在相同条件下重复进行；
	\item 试验结果可能不止一个，但能明确所有的可能结果；
	\item 试验前无法确定是哪个结果会出现。
\end{enumerate}

\begin{mydefinition}{样本空间与样本点}
	试验的所有可能结果所组成的集合，称为\textbf{样本空间}，记为
	\[\Omega = \{\omega\}.\]
	其中每个元素 $\omega$ 称为\textbf{样本点}。
\end{mydefinition}

\begin{mydefinition}{随机事件}
	试验中可能出现也可能不出现的情况称为随机事件，记作 $A,B,C$ 等。任何事件均对应着样本空间的某个\textbf{子集}。
\end{mydefinition}

\textbf{特殊事件：}
\begin{itemize}
	\item \textbf{必然事件：} 样本空间 $\Omega$；
	\item \textbf{不可能事件：} 空集 $\varnothing$；
	\item \textbf{基本事件：} 只包含一个样本点的事件 $\{\omega\}$。
\end{itemize}

\begin{example}[掷骰子]
掷一颗骰子，考察可能出现的点数。
\begin{itemize}
	\item 样本空间：$S_4 = \{1,2,3,4,5,6\}$；
	\item 事件 $A = \text{``掷出偶数点''} = \{2,4,6\}$；
	\item 事件 $B = \text{``掷出大于4的点''} = \{5,6\}$。
\end{itemize}
\end{example}

\section{事件间的关系与运算}

\subsection{基本关系}

\begin{enumerate}
	\item \textbf{包含关系：} $A \subset B$（$A$ 发生必导致 $B$ 发生）；
	\item \textbf{和事件：} $A \cup B$，$A$ 与 $B$ 至少有一个发生；
	\item \textbf{积事件：} $AB = A \cap B$，$A$ 与 $B$ 同时发生；
	\item \textbf{差事件：} $A - B$，$A$ 发生而 $B$ 不发生；
	\item \textbf{互斥事件：} $AB = \varnothing$；
	\item \textbf{互逆事件：} $A \cup B = \Omega$ 且 $AB = \varnothing$，此事 $A$ 与 $B$ 互为逆事件，记为 $B = \overline{A}$。
\end{enumerate}

\subsection{事件的运算律}

\begin{myproperty}{事件运算法则}
\begin{enumerate}
	\item \textbf{交换律：} $A \cup B = B \cup A$，$AB = BA$；
	\item \textbf{结合律：} $(A \cup B) \cup C = A \cup (B \cup C)$，$(AB)C = A(BC)$；
	\item \textbf{分配律：} $(A \cup B)C = (AC) \cup (BC)$，$(AB) \cup C = (A \cup C)(B \cup C)$；
	\item \textbf{对偶律 (De~Morgan 定律)：}
	\[\overline{A \cup B} = \overline{A} \cap \overline{B},\qquad \overline{AB} = \overline{A} \cup \overline{B}\]
	可推广为：
	\[\overline{\bigcup_{k} A_k} = \bigcap_{k} \overline{A}_k,\qquad
	\overline{\bigcap_{k} A_k} = \bigcup_{k} \overline{A}_k\]
\end{enumerate}
\end{myproperty}

\section{概率的定义与性质}

\subsection{Kolmogorov 公理化定义}

设 $\Omega$ 为样本空间，$\mathcal{F}$ 为事件域（$\sigma$-代数），概率 $P$ 是定义在 $\mathcal{F}$ 上的实值函数，满足：

\begin{myhypothesis}{概率公理（Kolmogorov 1933）}
\begin{enumerate}
	\item \textbf{非负性：} $\forall A \in \mathcal{F}$，$P(A) \geq 0$；
	\item \textbf{规范性：} $P(\Omega) = 1$；
	\item \textbf{可列可加性：} 若 $A_1, A_2, \ldots$ 两两互不相容，则
	\[P\!\left(\bigcup_{i=1}^{\infty} A_i\right) = \sum_{i=1}^{\infty} P(A_i).\]
\end{enumerate}
\end{myhypothesis}

\subsection{概率的基本性质}

\begin{myproperty}{概率性质}
\begin{enumerate}
	\item $P(\varnothing) = 0$；
	\item \textbf{有限可加性：} 若 $A_1, \ldots, A_n$ 两两互不相容，则
	\[P(A_1 \cup A_2 \cup \cdots \cup A_n) = \sum_{i=1}^{n} P(A_i);\]
	\item \textbf{事件差公式：} $P(A - B) = P(A) - P(AB)$；
	\item \textbf{加法公式：} 对任意两事件 $A, B$，有
	\[P(A \cup B) = P(A) + P(B) - P(AB);\]
	推广到三个事件：
	\begin{align*}
	P(A \cup B \cup C) &= P(A) + P(B) + P(C) \\
	&\quad - P(AB) - P(AC) - P(BC) + P(ABC);
	\end{align*}
	\item \textbf{互补性：} $P(\overline{A}) = 1 - P(A)$。
\end{enumerate}
\end{myproperty}

\section{古典概型}

若样本空间 $\Omega$ 只包含有限个样本点，且每个样本点出现的可能性相等，则称为\textbf{古典概型}。事件 $A$ 的概率为：
\begin{myformula}
	\boxed{\displaystyle P(A) = \frac{|A|}{|\Omega|} = \frac{\text{$A$ 包含的样本点数}}{\text{样本点总数}}}
\end{myformula}

\begin{example}[古典概型抽球公式]
盒中有 $N$ 个球，其中 $M$ 个白球。现从中任抽 $n$ 个球，则这 $n$ 个球中恰有 $k$ 个白球的概率为：
\[P = \frac{C_M^k \cdot C_{N-M}^{n-k}}{C_N^n}.\]
\end{example}

\section{条件概率与独立性}

\begin{mydefinition}{条件概率}
	设 $A, B$ 为两个事件，且 $P(A) > 0$，则称
	\[P(B \mid A) = \frac{P(AB)}{P(A)}\]
	为在事件 $A$ 发生的条件下事件 $B$ 发生的条件概率。
\end{mydefinition}

\begin{theorem}[乘法公式与链式法则]
由条件概率定义直接可得：
\[P(AB) = P(A) \cdot P(B \mid A).\]
一般形式（链式法则）：
\[P(A_1 A_2 \cdots A_n) = P(A_1) P(A_2 \mid A_1) P(A_3 \mid A_1 A_2) \cdots P(A_n \mid A_1 \cdots A_{n-1}).\]
\end{theorem}

\begin{mydefinition}{事件的独立性}
	若 $P(AB) = P(A) P(B)$，则称事件 $A$ 与 $B$ 相互独立，记为 $A \indep B$。
\end{mydefinition}

\begin{mydefinition}{多个事件的独立性}
	设 $A_1, A_2, \ldots, A_n$ 为 $n$ 个事件。若对任意 $k$ ($2 \leq k \leq n$) 和任意 $1 \leq i_1 < i_2 < \cdots < i_k \leq n$，有
	\[P(A_{i_1} A_{i_2} \cdots A_{i_k}) = P(A_{i_1}) P(A_{i_2}) \cdots P(A_{i_k}),\]
	则称 $A_1, A_2, \ldots, A_n$ \textbf{相互独立}。
\end{mydefinition}

\begin{remark}
	两两独立不等价于相互独立。相互独立要求任意子集的乘积概率等于各自概率的乘积。
\end{remark}

\begin{example}[独立性证明]
证明：若 $A, B$ 相互独立，则 $A$ 与 $\overline{B}$ 相互独立。

\begin{solution}
\begin{align*}
P(A\overline{B}) &= P(A - AB) = P(A) - P(AB) \qquad\text{(事件差公式)}\\
&= P(A) - P(A)P(B) \qquad\text{(因 $A, B$ 独立)}\\
&= P(A)[1 - P(B)] = P(A)P(\overline{B}).
\end{align*}
故 $A$ 与 $\overline{B}$ 相互独立。
\end{solution}
\end{example}

\section{全概率公式与贝叶斯公式}

\begin{theorem}[全概率公式]
设 $A_1, A_2, \ldots, A_n$ 为样本空间 $\Omega$ 的一个完备事件组（即两两互斥且 $\bigcup_{i=1}^{n} A_i = \Omega$），且 $P(A_i) > 0$，则对任意事件 $B$，有
\[P(B) = \sum_{i=1}^{n} P(A_i) P(B \mid A_i).\]
\end{theorem}

\begin{theorem}[贝叶斯公式]
在全概率公式的相同条件下，有
\[P(A_j \mid B) = \frac{P(A_j) P(B \mid A_j)}{\displaystyle\sum_{i=1}^{n} P(A_i) P(B \mid A_i)},\quad j = 1, 2, \ldots, n.\]
\end{theorem}

\begin{myTheorem}{贝叶斯公式的理解}
\begin{itemize}
	\item $P(A_j)$ 称为\textbf{先验概率}（prior），是试验前对 $A_j$ 的认知；
	\item $P(B \mid A_j)$ 称为\textbf{似然}（likelihood），是新数据带来的信息；
	\item $P(A_j \mid B)$ 称为\textbf{后验概率}（posterior），是在观察到 $B$ 发生之后对 $A_j$ 的更新判断。
\end{itemize}
贝叶斯公式本质上是一种"用新数据更新旧信念"的数学机制。
\end{myTheorem}

\begin{example}[袋子抽球——全概率公式]
甲袋（2白，1红），乙袋（2红，1白）。从甲袋中任取一球放入乙袋，再从乙袋任取一球是红球（$B$）的概率。

设 $A_1$——从甲袋放入乙袋的是白球，$P(A_1) = 2/3$；$A_2$——从甲袋放入乙袋的是红球，$P(A_2) = 1/3$。

\begin{itemize}
	\item 若 $A_1$ 发生（乙袋变为2白2红）：$P(B \mid A_1) = 2/4 = 1/2$；
	\item 若 $A_2$ 发生（乙袋变为1白3红）：$P(B \mid A_2) = 3/4$。
\end{itemize}

由全概率公式：
\[P(B) = P(B \mid A_1)P(A_1) + P(B \mid A_2)P(A_2) = \frac{1}{2} \cdot \frac{2}{3} + \frac{3}{4} \cdot \frac{1}{3} = \frac{7}{12}.\]
\end{example}

\begin{remark}[概率为0的事件]
概率为 $0$ 的事件不一定是不可能事件。例如在连续型随机变量中，取任意指定值的概率均为 $0$，但该事件仍可能发生。
\end{remark}
